\section{Marco teórico y trabajos relacionados}
Sistemas de información para transporte. Distintos proyectos han mostrado beneficios de publicar posiciones en tiempo real y permitir que el pasajero visualice rutas y buses activos \cite{cityruta2018new}; \cite{stopbus2016new}. Estas experiencias resaltan mejoras en la percepción del servicio y en la planificación del viaje.

Mensajería ligera y pub/sub. En contextos móviles y de IoT, usar MQTT con un broker como Mosquitto permite distribuir eventos con baja sobrecarga y buena escalabilidad. Trabajos de seguridad y rastreo vehicular evidencian que el patrón publish/subscribe es adecuado para manejar muchas actualizaciones de ubicación \cite{sismo2016}. Esta idea también se usa en módulos de comunicación genéricos para aplicaciones conectadas y en soluciones académicas/industriales.

APIs claras y mantenibilidad. Documentar el contrato de la API (OpenAPI/Swagger) reduce ambigüedades y acelera la integración \cite{easyrest2022}. La literatura sobre diseño orientado a objetos y sostenibilidad del software pone énfasis en separar responsabilidades y favorecer la extensibilidad \cite{poo2023}.

Arquitecturas híbridas y datos. Decidir cuándo usar almacenamiento relacional y cuándo incorporar componentes NoSQL es clave al escalar servicios de rastreo. Guías sobre sistemas híbridos ayudan a elegir la mezcla adecuada según el tipo de consulta y el volumen de eventos \cite{ingenieria2022}.

Escalabilidad en la práctica. Migraciones exitosas de monolitos a microservicios en sistemas de transporte (como el caso Kbus) reportan millones de registros diarios y cientos de miles de consultas, lo que prueba que estos dominios pueden crecer mucho \cite{torres2020}. En logística, soluciones cloud de tracking confirman el valor de ver el momento exacto de los movimientos y notificar a usuarios clave \cite{odesanya2025}. También hay experiencias recientes de rastreo en tiempo real en contextos con infraestructura desafiante \cite{anyango2025}.

Síntesis y vacíos. En conjunto, la literatura respalda: (i) publicar datos en vivo; (ii) usar canales ligeros (MQTT/WebSockets); (iii) documentar APIs; y (iv) diseñar para escalar. Persisten vacíos en guías paso a paso que integren todas estas piezas para contextos locales con presupuesto limitado. UrbanTracker llena parte de ese espacio con una receta técnica concreta y validaciones iniciales.